\documentclass[11pt]{article}
\usepackage[a4paper,margin=1in]{geometry}
\usepackage{fontspec}
\usepackage[boldfont=false]{xeCJK}
\setCJKmainfont{FandolSong-Regular.otf}
\usepackage{amsmath}
\usepackage{amssymb}
\usepackage{array}
\usepackage{booktabs}
\usepackage{graphicx}
\usepackage{adjustbox}
\usepackage{enumitem}
\setlist[itemize]{topsep=2pt,partopsep=0pt,parsep=0pt,itemsep=2pt,leftmargin=1.4em}
\setlist[enumerate]{topsep=2pt,partopsep=0pt,parsep=0pt,itemsep=2pt,leftmargin=1.6em}
\usepackage{parskip}
\usepackage{fvextra}
\fvset{breaklines=true,breakanywhere=true,fontsize=\small}
\setlength{\parindent}{0pt}

\begin{document}

\begin{center}
  {\LARGE\bfseries Group 2}\\[0.35em]
  {\large A Level Chemistry}
\end{center}
\vspace{0.6em}

\section*{Thermal stability of the nitrates and carbonates}

The \textbf{thermal stability} 热稳定性 of the Group 2 \textbf{nitrates} 硝酸盐 and \textbf{carbonates} 碳酸盐 \textbf{increases} down the group. Here is the reason, in terms of how the ions affect each other.

A small \textbf{cation} 阳离子 with a small \textbf{ionic radius} 离子半径 has a high charge density. It pulls on the electrons of the nearby anion and distorts its shape — this is \textbf{polarisation} 极化. Distorting the large \textbf{anion} 阴离子 (the carbonate or nitrate ion) weakens a bond inside it, so the compound breaks down more easily.

Going down the group, the cation gets larger. Its charge density drops, so it polarises the anion \textbf{less}. The anion is less distorted, so the compound is harder to break down — it is more thermally stable and needs a higher temperature to decompose.

\section*{Solubility of the hydroxides and sulfates}

When an ionic solid dissolves, two energy changes compete. The \textbf{enthalpy change of solution} 溶解焓变 is the sum of them:

\[
\Delta H^{\ominus}_{\text{sol}} = -\Delta H_{\text{latt}} + \Delta H_{\text{hyd}}
\]

\begin{itemize}
\item you must first pull the lattice apart (this needs the \textbf{lattice energy} 晶格能).

\item then water surrounds the ions (this releases the \textbf{enthalpy change of hydration} 水合焓变).

\end{itemize}

If the energy released on hydration roughly matches (or beats) the energy needed to break the lattice, the solid dissolves easily.

Going down the group, both the lattice energy and the hydration enthalpy get \textbf{smaller} (less negative), because the cation is larger. But they shrink at different rates, and this explains the opposite trends:

\begin{center}
\adjustbox{max width=\linewidth}{%
\begin{tabular}{lll}
\toprule
\textbf{Compound} & \textbf{Trend in \textbf{solubility} 溶解度 down the group} & \textbf{Why} \\
\midrule
\textbf{hydroxides} 氢氧化物 & increases & the $\text{OH}^-$ ion is small, so the lattice energy falls a lot down the group; this change outweighs the fall in hydration energy \\
\textbf{sulfates} 硫酸盐 & decreases & the $\text{SO}_4^{2-}$ ion is large, so the lattice energy stays almost the same; the fall in hydration energy then dominates, so dissolving becomes less favourable \\
\bottomrule
\end{tabular}}
\end{center}


\end{document}
